% engines: lualatex
% expect-warnings: 1
% expect-errors: 0
\documentclass{article}
\usepackage[margin=1cm]{geometry}
\usepackage[kernel = lua]{../../latex/pgf-go}
\newcommand\case[1]{\par\medskip\noindent\textbf{#1}\par\nopagebreak}
\pagestyle{empty}
\begin{document}

% game.sgf: 9x9; B E5, W G7, B pass, W C3, B G5.
% ko.sgf: 5x5 ko; W C3, B D3 (captures 1), W C3 (retakes, captures 2), B A1, W E5.

\case{game.sgf, step = 2: diagrams 1--2, 3--4 (3 is a pass) and 5; the earlier moves have no numbers}
\sgfdiagrams[step = 2]{game.sgf}

\case{step = 2 and restart: each diagram starts at 1 (4 is shown as 2, 5 as 1) / to = 3: two diagrams}
\sgfdiagrams[step = 2, restart]{game.sgf}

\sgfdiagrams[step = 2, to = 3]{game.sgf}

\case{ko.sgf in one diagram: 2 is not removed, 3 is noted ("3 en 1"); title = Diagrama}
\begingroup
\pgfgoset{notes = {at = en}}
\sgfdiagrams[step = 5, title = Diagrama]{ko.sgf}
\endgroup

\case{ko.sgf, step = 2: diagram 2 starts without 1 (captured), and 2 is kept when 3 retakes; notes = false}
\sgfdiagrams[step = 2, notes = false]{ko.sgf}

\case{goban options for every diagram, and caption font}
\sgfdiagrams[step = 5, goban = {label = {at = {north, west}}}, caption = {font = \itshape}]{ko.sgf}

\case{step = \{1, 3\} and restart = \{from = 0\}: 1; 2--4 (2 is shown as 0, 4 as 2); 5 (shown as 0)}
\sgfdiagrams[step = {1, 3}, restart = {from = 0}]{game.sgf}

\case{step = 2, from = 4, restart = \{from = 0\}: diagram 2 (4--4), 4 is shown as 1; diagram 3, 5 as 0}
\sgfdiagrams[step = 2, from = 4, restart = {from = 0}]{game.sgf}

\case{restart = \{from = 0\}, then restart = false: no restart (as the first diagram)}
\sgfdiagrams[step = 2, restart = {from = 0}, restart = false]{game.sgf}

\newcommand\notes[1][]{\par\noindent Notes: [\pgfgomovenotes[#1]]\par}
\case{goban/notes with a pass: "3: pass" / from = 4: not noted / numbers = false: no notes / pass = pasa: "3: pasa"}
\begin{goban}[9, notes]\sgf{game.sgf}\end{goban}
\notes
\begin{goban}[9, notes]\sgf[from = 4]{game.sgf}\end{goban}
\notes
\begin{goban}[9, notes]\sgf[numbers = false]{game.sgf}\end{goban}
\notes
\begin{goban}[9, notes]\sgf{game.sgf}\end{goban}
\notes[pass = pasa]

% var-tree.sgf: 9x9; 1 B E5, 2 W C7, then three lines: main 3 B G3, 4 W G7 "Main.",
% 5 B C3; variation 2: 3 B G7 "Variation 2, commented.", 4 W G3, 5 B C3, with the
% sub-variation 2.2: 4 W C3 "Sub-variation 2.2.", 5 B G3; variation 3: 3 B C3, 4 W G3.

\case{variations (the commented ones, depth 1): after diagram 2, diagram 2a (variation 2, numbered from 1, "variation at 3")}
\sgfdiagrams[step = 2, comments, variations, goban = {scale = 0.5}]{var-tree.sgf}

\case{variations = \{all, depth = 2\}: 2a (2), 2b (2.2, from its move 4: "1: Sub-variation"), 2c (3)}
\sgfdiagrams[step = 2, comments, variations = {all, depth = 2}, goban = {scale = 0.5}]{var-tree.sgf}

\case{variations = \{3, 2.2.1\}: 2a (2.2), 2b (3), in the order of the tree / \{7\}: one warning, no variation}
\sgfdiagrams[step = 2, variations = {3, 2.2.1}, goban = {scale = 0.5}]{var-tree.sgf}
\sgfdiagrams[step = 5, variations = {7}, goban = {scale = 0.5}]{var-tree.sgf}

\end{document}
